% GENERATED FILE. Edit canonical lessons, then run npm run generate:books.
% canonical-source-hash: fnv1a64:15f0ca0f28ee3cc6
% canonical-lessons: SA-C359-tyagah, SA-C359-paropakarah, SA-C359-svarthah, SA-C359-trsna, SA-C359-samarpayati

\chapter{Giving Up for Others}
\label{ch:sa-giving-up-for-others}
\hlchaptermodality{359}

\begin{chapteropening}
\textbf{By the end of this chapter:} \emph{I can talk about selfless and selfish acts.}

Complete the last lesson of chapter 359: I can talk about selfless and selfish acts.
\end{chapteropening}

\section[\texorpdfstring{\sk{त्यागः} (tyāgaḥ)}{त्यागः (tyāgaḥ)}]{\sk{त्यागः} (tyāgaḥ) — renunciation; giving up}

\label{lesson:SA-C359-tyagah}



\begin{quote}
Before the new one: say the Sanskrit for the right path, then the Sanskrit for holds in regard.
\end{quote}

\subsection*{You'll want to know: \sk{त्यागः}}
\textbf{\sk{त्यागः}} — \emph{tyāgaḥ} — \textquotedblleft{}renunciation; giving up\textquotedblright{}.

\begin{quote}
\sk{सैनिकस्य} \sk{त्यागः} \sk{महान्} \sk{भवति।}
\end{quote}

\begin{quote}
\emph{sainikasya tyāgaḥ mahān bhavati.}
\end{quote}

\begin{quote}
A soldier's sacrifice is great.
\end{quote}

A masculine noun.

\begin{cousinweb}
\textbf{\sk{त्यागः}} (\emph{tyāgaḥ}) comes from the root √tyaj, \textquotedblleft{}to leave, to abandon\textquotedblright{}; MW gives \emph{tyāga} \textquotedblleft{}leaving, abandoning\textquotedblright{} and \textquotedblleft{}giving up, a gift\textquotedblright{} (MW).
\end{cousinweb}

\begin{grammarlens}[title={What you've built}]
One more word for giving reasons and opinions.
\end{grammarlens}

\subsection*{Your turn}
\begin{itemize}
\raggedright
  \item \emph{Say it:} \emph{tyāgaḥ}
  \item \emph{Say it:} \emph{tyāgaḥ}, once more
  \item \emph{Say it:} \emph{sainikasya tyāgaḥ mahān bhavati.} — A soldier's sacrifice is great.
  \item \emph{Recall:} say the Sanskrit for the right path, then the Sanskrit for holds in regard, then say \emph{tyāgaḥ} again
\end{itemize}

\begin{tcolorbox}[breakable,colback=teal!4,colframe=teal!35!black,title={Before you move on}]
What does \sk{त्यागः} mean? (\textquotedblleft{}Renunciation; giving up\textquotedblright{}.) Say it once more.
\end{tcolorbox}
\section[\texorpdfstring{\sk{परोपकारः} (paropakāraḥ)}{परोपकारः (paropakāraḥ)}]{\sk{परोपकारः} (paropakāraḥ) — helping others}

\label{lesson:SA-C359-paropakarah}



\begin{quote}
Before the new one: say the Sanskrit for holds in regard, then the Sanskrit for renunciation.
\end{quote}

\subsection*{You'll want to know: \sk{परोपकारः}}
\textbf{\sk{परोपकारः}} — \emph{paropakāraḥ} — \textquotedblleft{}helping others\textquotedblright{}.

\begin{quote}
\sk{परोपकारः} \sk{पुण्यं} \sk{भवति।}
\end{quote}

\begin{quote}
\emph{paropakāraḥ puṇyaṃ bhavati.}
\end{quote}

\begin{quote}
Helping others earns merit.
\end{quote}

A masculine noun.

\begin{cousinweb}
\textbf{\sk{परोपकारः}} (\emph{paropakāraḥ}) is built from \emph{para}, \textquotedblleft{}another\textquotedblright{}, and \emph{upakāra}, \textquotedblleft{}a favour, help\textquotedblright{}; MW gives \emph{paropakāra} \textquotedblleft{}assisting others, benevolence, charity\textquotedblright{} (MW).
\end{cousinweb}

\begin{grammarlens}[title={What you've built}]
One more word for giving reasons and opinions.
\end{grammarlens}

\subsection*{Your turn}
\begin{itemize}
\raggedright
  \item \emph{Say it:} \emph{paropakāraḥ}
  \item \emph{Say it:} \emph{paropakāraḥ}, once more
  \item \emph{Say it:} \emph{paropakāraḥ puṇyaṃ bhavati.} — Helping others earns merit.
  \item \emph{Recall:} say the Sanskrit for holds in regard, then the Sanskrit for renunciation, then say \emph{paropakāraḥ} again
\end{itemize}

\begin{tcolorbox}[breakable,colback=teal!4,colframe=teal!35!black,title={Before you move on}]
What does \sk{परोपकारः} mean? (\textquotedblleft{}Helping others\textquotedblright{}.) Say it once more.
\end{tcolorbox}
\section[\texorpdfstring{\sk{स्वार्थः} (svārthaḥ)}{स्वार्थः (svārthaḥ)}]{\sk{स्वार्थः} (svārthaḥ) — self-interest}

\label{lesson:SA-C359-svarthah}



\begin{quote}
Before the new one: say the Sanskrit for renunciation, then the Sanskrit for helping others.
\end{quote}

\subsection*{You'll want to know: \sk{स्वार्थः}}
\textbf{\sk{स्वार्थः}} — \emph{svārthaḥ} — \textquotedblleft{}self-interest\textquotedblright{}.

\begin{quote}
\sk{सः} \sk{केवलं} \sk{स्वार्थं} \sk{चिन्तयति।}
\end{quote}

\begin{quote}
\emph{saḥ kevalaṃ svārthaṃ cintayati.}
\end{quote}

\begin{quote}
He thinks only of his own interest.
\end{quote}

A masculine noun.

\begin{cousinweb}
\textbf{\sk{स्वार्थः}} (\emph{svārthaḥ}) is built from \emph{sva}, \textquotedblleft{}own\textquotedblright{}, and \emph{artha}, \textquotedblleft{}aim, advantage\textquotedblright{}; MW gives \emph{svārtha} \textquotedblleft{}one's own affair, personal advantage, self-interest\textquotedblright{} (MW).
\end{cousinweb}

\begin{grammarlens}[title={What you've built}]
One more word for giving reasons and opinions.
\end{grammarlens}

\subsection*{Your turn}
\begin{itemize}
\raggedright
  \item \emph{Say it:} \emph{svārthaḥ}
  \item \emph{Say it:} \emph{svārthaḥ}, once more
  \item \emph{Say it:} \emph{saḥ kevalaṃ svārthaṃ cintayati.} — He thinks only of his own interest.
  \item \emph{Recall:} say the Sanskrit for renunciation, then the Sanskrit for helping others, then say \emph{svārthaḥ} again
\end{itemize}

\begin{tcolorbox}[breakable,colback=teal!4,colframe=teal!35!black,title={Before you move on}]
What does \sk{स्वार्थः} mean? (\textquotedblleft{}Self-interest\textquotedblright{}.) Say it once more.
\end{tcolorbox}
\section[\texorpdfstring{\sk{तृष्णा} (tṛṣṇā)}{तृष्णा (tṛṣṇā)}]{\sk{तृष्णा} (tṛṣṇā) — craving}

\label{lesson:SA-C359-trsna}



\begin{quote}
Before the new one: say the Sanskrit for helping others, then the Sanskrit for self-interest.
\end{quote}

\subsection*{You'll want to know: \sk{तृष्णा}}
\textbf{\sk{तृष्णा}} — \emph{tṛṣṇā} — \textquotedblleft{}craving\textquotedblright{}.

\begin{quote}
\sk{तृष्णा} \sk{दुःखस्य} \sk{मूलम्} \sk{अस्ति।}
\end{quote}

\begin{quote}
\emph{tṛṣṇā duḥkhasya mūlam asti.}
\end{quote}

\begin{quote}
Craving is the root of sorrow.
\end{quote}

A feminine noun.

\begin{cousinweb}
\textbf{\sk{तृष्णा}} (\emph{tṛṣṇā}) comes from the root √tṛṣ, \textquotedblleft{}to be thirsty\textquotedblright{}; MW gives \emph{tṛṣṇā} \textquotedblleft{}thirst\textquotedblright{} and \textquotedblleft{}desire, avidity\textquotedblright{} (MW).
\end{cousinweb}

\begin{grammarlens}[title={What you've built}]
One more word for giving reasons and opinions.
\end{grammarlens}

\subsection*{Your turn}
\begin{itemize}
\raggedright
  \item \emph{Say it:} \emph{tṛṣṇā}
  \item \emph{Say it:} \emph{tṛṣṇā}, once more
  \item \emph{Say it:} \emph{tṛṣṇā duḥkhasya mūlam asti.} — Craving is the root of sorrow.
  \item \emph{Recall:} say the Sanskrit for helping others, then the Sanskrit for self-interest, then say \emph{tṛṣṇā} again
\end{itemize}

\begin{tcolorbox}[breakable,colback=teal!4,colframe=teal!35!black,title={Before you move on}]
What does \sk{तृष्णा} mean? (\textquotedblleft{}Craving\textquotedblright{}.) Say it once more.
\end{tcolorbox}
\section[\texorpdfstring{\sk{समर्पयति} (samarpayati)}{समर्पयति (samarpayati)}]{\sk{समर्पयति} (samarpayati) — dedicates; hands over}

\label{lesson:SA-C359-samarpayati}



\begin{quote}
Before the new one: say the Sanskrit for self-interest, then the Sanskrit for craving.
\end{quote}

\subsection*{You'll want to know: \sk{समर्पयति}}
\textbf{\sk{समर्पयति}} — \emph{samarpayati} — \textquotedblleft{}dedicates; hands over\textquotedblright{}.

\begin{quote}
\sk{सः} \sk{धनं} \sk{विद्यालयाय} \sk{समर्पयति।}
\end{quote}

\begin{quote}
\emph{saḥ dhanaṃ vidyālayāya samarpayati.}
\end{quote}

\begin{quote}
He hands over the money to the school.
\end{quote}

Past (he …): \sk{सः} \sk{गृहं} \sk{समर्पितवान्}, \emph{saḥ gṛhaṃ samarpitavān}.

\begin{cousinweb}
\textbf{\sk{समर्पयति}} (\emph{samarpayati}) is built on \emph{sam} and the root √ṛ, \textquotedblleft{}to go\textquotedblright{}, in its causative form \emph{arpaya}, \textquotedblleft{}to hand over\textquotedblright{}; MW gives the sense \textquotedblleft{}to deliver over, to consign, to commit\textquotedblright{} (MW).
\end{cousinweb}

\begin{grammarlens}[title={What you've built}]
One more word for giving reasons and opinions.
\end{grammarlens}

\subsection*{Your turn}
\begin{itemize}
\raggedright
  \item \emph{Say it:} \emph{samarpayati}
  \item \emph{Say it:} \emph{samarpayati}, once more
  \item \emph{Say it:} \emph{saḥ dhanaṃ vidyālayāya samarpayati.} — He hands over the money to the school.
  \item \emph{Recall:} say the Sanskrit for self-interest, then the Sanskrit for craving, then say \emph{samarpayati} again
\end{itemize}

\begin{tcolorbox}[breakable,colback=teal!4,colframe=teal!35!black,title={Before you move on}]
What does \sk{समर्पयति} mean? (\textquotedblleft{}Dedicates; hands over\textquotedblright{}.) Say it once more.
\end{tcolorbox}
